\section{Metodologia de validação}
\label{sec:debate-valid}

\begin{confronto}
Como o projeto valida os modelos? Há um conjunto de teste honesto? A separação
respeita a natureza temporal dos dados, ou há vazamento?
\end{confronto}

\subsection{Três níveis de separação}
A validação opera em três níveis, deliberadamente distintos:

\begin{enumerate}[noitemsep]
  \item \textbf{Treino} (2010--2019): aprende a normalidade.
  \item \textbf{Validação} (bloco final do treino, $\approx$2018--2019, via
        \texttt{validation\_split=0,1} com \texttt{shuffle=False}): governa o
        \emph{early stopping} e a calibração de hiperparâmetros. \textbf{Não} é o
        teste.
  \item \textbf{Teste} (2020--2024): nunca visto no treino; é onde as anomalias
        são detectadas e as métricas reportadas.
\end{enumerate}

\paragraph{O detalhe crítico: \texttt{shuffle=False}.}
O \texttt{validation\_split} do Keras separa a \emph{última fração} do array
\emph{antes} de embaralhar. Com \texttt{shuffle=True}, cada época re-embaralharia
o treino e --- combinada com janelas sobrepostas --- misturaria amostras
temporalmente adjacentes entre treino e validação, vazando padrão futuro. Com
\texttt{shuffle=False}, treino e validação ficam em blocos cronológicos
contíguos. Esta é a defesa central contra vazamento, e está registrada como
decisão crítica (ADR-0004).

\begin{honesto}
Há um vazamento \emph{residual} conhecido: janelas na fronteira treino/validação
compartilham alguns dias, por causa da sobreposição de janelas. O efeito é
pequeno e foi assumido e documentado, não ignorado.
\end{honesto}

\subsection{Validação cruzada para seleção de modelos}
\label{sec:debate-valid-cv}

\begin{confronto}
Para escolher o modelo ótimo entre opções pré-determinadas, usou-se
\emph{validação cruzada}? Onde está o $k$-fold?
\end{confronto}

Aqui é preciso ser direto: \textbf{não usamos $k$-fold clássico --- e isso é
intencional, não um esquecimento.}

\paragraph{Por que $k$-fold padrão é inválido aqui.}
A validação cruzada $k$-fold embaralha e particiona os dados assumindo que as
amostras são \emph{independentes e identicamente distribuídas} (i.i.d.). Séries
temporais violam isso: as observações são dependentes e ordenadas. Um \emph{fold}
de treino contendo dias \emph{posteriores} aos do \emph{fold} de teste constitui
vazamento de futuro --- exatamente o pecado da Seção~\ref{sec:debate-valid}. Em
detecção de anomalias isso é ainda pior, pois infla artificialmente a aparência
de desempenho.

\paragraph{O que de fato fizemos: seleção por \emph{grid} sobre validação temporal.}
A ``comparação de múltiplos modelos para escolher o ótimo'' aconteceu, mas pela
via apropriada a séries temporais --- avaliar candidatos numa validação que
respeita o tempo:
\begin{itemize}[noitemsep]
  \item \textbf{Dimensão latente} $\in\{8,16,32\}$: comparadas pela
        \texttt{val\_loss} no bloco de validação cronológico; resultado plano,
        16 confirmado (Seção~\ref{sec:debate-arq-prof}).
  \item \textbf{Esquema de limiar} (estático p95 vs.\ dinâmico causal) e
        \textbf{tamanho da janela dinâmica} $\in\{126,252,504\}$: comparados pela
        taxa de marcação e pelo comportamento sob mudança de regime no teste; 252
        confirmado.
  \item \textbf{Modelo profundo vs.\ \emph{baseline}}: o LSTM-AE é confrontado
        com um \emph{baseline} de $z$-score em janela móvel, para evidenciar
        ganho sobre a heurística trivial.
\end{itemize}

\paragraph{A forma correta de CV temporal (e nosso caminho de rigor).}
A validação cruzada que \emph{seria} válida aqui é a de série temporal ---
\emph{walk-forward} / janela expansível, em que cada \emph{fold} treina no
passado e valida no futuro imediato, nunca o contrário:

\begin{lstlisting}[caption={Validação cruzada temporal (walk-forward) com scikit-learn.}]
from sklearn.model_selection import TimeSeriesSplit
tscv = TimeSeriesSplit(n_splits=5)
for treino_idx, val_idx in tscv.split(X):
    # treino_idx sempre ANTES de val_idx -> sem vazamento de futuro
    modelo = build_lstm_autoencoder()
    modelo.fit(X[treino_idx], X[treino_idx], shuffle=False, ...)
    avalia(modelo, X[val_idx])
\end{lstlisting}

\begin{honesto}
O projeto usou um \textbf{único} \emph{holdout} temporal de validação (não uma
$k$-fold temporal com vários cortes). Para a seleção de hiperparâmetros conduzida
--- de baixa dimensionalidade e com resultados planos/robustos --- isso é
suficiente e defensável. Porém, um \texttt{TimeSeriesSplit} (\emph{walk-forward})
daria estimativas com \emph{variância} (média $\pm$ desvio entre cortes) e uma
seleção de modelo estatisticamente mais sólida. Está registrado como o próximo
passo de rigor --- e o ADR-0004 já o aponta como alternativa considerada.
\end{honesto}

\subsection{Validação substantiva: além da perda}
Vale frisar que a \texttt{val\_loss} não é o único juiz. Como não há rótulos de
anomalia reais, a validação do \emph{detector} (não só do reconstrutor) se dá por
duas vias complementares (detalhadas no relatório principal):
\begin{itemize}[noitemsep]
  \item \textbf{Quantitativa:} injeção sintética de anomalias em posições
        conhecidas $\to$ Precision/Recall/F1.
  \item \textbf{Narrativa:} as anomalias detectadas coincidem com eventos reais
        (COVID, fraude da Americanas, Brumadinho)?
\end{itemize}
Essa dupla checagem é a resposta honesta para ``como validar um modelo sem
gabarito''.
